% ==============================================================================
% چکیده فارسی
% ==============================================================================

\newpage
\thispagestyle{plain}

\begin{center}
    {\titr\Large چکیده}\\[1cm]
\end{center}

\noindent
در سامانه‌های تعبیه‌شده و بی‌درنگ سخت (\lr{Hard Real-Time Embedded Systems})، تخمین دقیق و تضمین‌پذیر بیشینه زمان اجرا (\lr{WCET}) عاملی حیاتی در اثبات ایمنی و زمان‌بندی‌پذیری سامانه به شمار می‌رود. در معماری‌های نوین پردازنده‌ها با خط‌لوله‌های چندمرحله‌ای، پیش‌بینی انشعاب و حافظه‌های نهان، روش‌های تحلیل ایستا به دلیل تخمین‌های بیش‌ازحد بدبینانه ناتوان بوده و در نقطه مقابل، رویکردهای سنتی ابزاردقیق نرم‌افزاری (\lr{Software Instrumentation}) با تحمیل سربار اجرایی و اعوجاج زمانی (پدیده «اثر پروب»)، رفتار اصیل پردازنده را دگرگون می‌سازند. این رساله به طراحی، پیاده‌سازی عملی و اعتبارسنجی جامع یک خط‌لوله غیرتهاجمی (\lr{Non-intrusive}) جهت تأمین داده‌های ردیابی اجرای برنامه و استنتاج دقیق \lr{WCET} در مقیاس چرخه ساعت پردازنده می‌پردازد.

سامانه ارائه‌شده با تمرکز بر ماژول ردیابی سخت‌افزاری \lr{ETMv4} در زیرساخت دیباگ \lr{ARM CoreSight} بر روی میکروکنترلر پیشرفته \lr{STM32H750 (ARM Cortex-M7)} توسعه یافته است. در این فرآیند، چالش‌های پیچیده پیکربندی ثبّاتی نظیر آزادسازی قفل‌های سخت‌افزاری، توزیع توان و کلاک و مغایرت‌های راهنمای مرجع تراشه با استاندارد معماری برطرف گردید. داده‌های فشرده ردیابی دستورات از طریق گذرگاه موازی ۴ بیتی درگاه \lr{TPIU} بدون تحمیل کوچک‌ترین سربار پردازشی صادر شده و توسط لاجیک‌آنالایزر دیجیتال با راهبرد اورسمپلینگ ۱۰ برابری ثبت می‌شوند. اسکریپت تخصصی \lr{\texttt{TraceStreamProcessor.py}} با ترازسازی فریم‌ها و جبران اعوجاج‌های فیزیکی، ساختار اسنپ‌شات استاندارد صنعتی را تولید کرده و با ادغام با موتور دکودر مرجع \lr{OpenCSD} و باینری \lr{ELF}، جریان تفکیک‌شده دستورات اسمبلی و برچسب‌های زمانی نانوثانیه‌ای را بازسازی می‌نماید. 

به منظور تبدیل این خط‌لوله به یک ابزار صنعتی و کاربرپسند، افزونه اختصاصی \lr{CoreSight Trace Studio} برای محیط توسعه \lr{VS Code} طراحی و پیاده‌سازی شد که در قالبی سه‌مرحله‌ای، متوالی و نیمه‌خودکار، فرآیند تزریق خودکار درایورها، نشانه‌گذاری پنجره‌های پایش در کدهای منبع، تولید متناظر مانیفست‌های اسنپ‌شات و نمایش لاگ‌های روند اجرای برنامه را در یک داشبورد مدرن محقق می‌سازد. در فاز ارزیابی تجربی، عملکرد سامانه تحت سناریوهای چالش‌برانگیز شامل آزمون دوره‌ای تایمر ۱۰۰ میکروثانیه با ثبت انحراف زمانی ناچیز زیر $0.1\ \mu\mathrm{s}$ (اثبات قطعی ماهیت غیرتهاجمی)، تطبیق زمانی سیگنال‌های فیزیکی، و رهگیری وقایع ترکیبی وقفه‌ها با هندشیک سخت‌افزاری به اثبات رسید. این بستر در هم‌افزایی با سامانه مکمل آزمون خودکار، پلتفرمی اقتصادی، بومی و کارا جهت آزمون برنامه‌های سیستم‌های نهفته و سنجش ایمنی آن‌ها به ارمغان آورده است. چه بسا این سیستم بتواند در آینده در قامت یک رقیب برای ابزارهای آزمون نرم‌افزار تجاری بین‌المللی واقع شود.

\vspace{1.0cm}
\noindent
\textbf{واژگان کلیدی:}\quad بیشینه زمان اجرا (\lr{WCET})، ردیابی غیرتهاجمی دستورات (\lr{Instruction Trace})، ماژول \lr{ETMv4}، معماری \lr{ARM CoreSight}، پردازنده \lr{Cortex-M7}، میکروکنترلر \lr{STM32H7}، دکودر \lr{OpenCSD}، افزونه \lr{CoreSight Trace Studio}، سامانه‌های بی‌درنگ سخت.
